The two target rows are linearly independent because, by \(\texttt{def:bow-target-coefficient}\),
\[
 r_0=\mathbf 1\otimes e_0,
 \qquad
 r_1=\mathbf 1\otimes e_1,
\]
and \(e_0,e_1\) are linearly independent.  Since both rows belong to \(\operatorname{row}(A_{\mathrm b})\), it follows that \(\operatorname{rank}(A_{\mathrm b})\ge2\).

Suppose for contradiction that \(\operatorname{rank}(A_{\mathrm b})=2\).  The two-dimensional inclusion then forces
\[
 \operatorname{row}(A_{\mathrm b})
 =\operatorname{span}\{r_0,r_1\}
 =\{\mathbf 1\otimes s:s\in\mathbb R^2\}.
\]
Fix a tagged treatment branch \(\beta_X\).  By \(\texttt{def:bow-probe-matrix}\), its row is the product vector
\[
 \phi_{X_{\mathrm b},\beta_X}\otimes\psi_{X_{\mathrm b},\beta_X}.
\]
The state \(\psi_{X_{\mathrm b},\beta_X}\) is a probability vector, so some coordinate \(z_*\in\{0,1\}\) has \(\psi_{X_{\mathrm b},\beta_X}(z_*)>0\).  Membership of the product row in \(\{\mathbf 1\otimes s:s\in\mathbb R^2\}\) says that its two treatment-output coordinates at this fixed \(z_*\) are equal.  Therefore
\[
 \phi_{X_{\mathrm b},\beta_X}(0)\psi_{X_{\mathrm b},\beta_X}(z_*)
 =\phi_{X_{\mathrm b},\beta_X}(1)\psi_{X_{\mathrm b},\beta_X}(z_*).
\]
Division by the strictly positive displayed state coordinate gives
\[
 \phi_{X_{\mathrm b},\beta_X}(0)=\phi_{X_{\mathrm b},\beta_X}(1),
\]
so \(\phi_{X_{\mathrm b},\beta_X}\) is proportional to \(\mathbf 1\).  This argument applies to every tagged branch, including a zero effect, which is the zero multiple of \(\mathbf 1\).  Hence all tagged treatment effects lie in \(\operatorname{span}\{\mathbf 1\}\), contradicting their assumed span \(\mathbb R^2\).  Thus rank two is impossible, and the already proved lower bound of two sharpens to
\[
 \operatorname{rank}(A_{\mathrm b})\ge3.
\]
Finally, \(A_{\mathrm b}\) has one row per tagged treatment branch, and the rank of a matrix is at most its number of rows.  Every normalized menu satisfying the hypotheses therefore has at least three tagged treatment branches.  Notice that normalization was not needed for the stronger rank implication itself; it only places the consequence in the normalized design class of interest.